\chapter*{Conclusion générale et Perspectives}
\addcontentsline{toc}{chapter}{Conclusion générale et Perspectives}

Ce projet de fin d'année a permis la conception et la réalisation de la plateforme PAGe, une solution intégrée de gouvernance algorithmique publique répondant au besoin d'évaluation, de suivi et d'atténuation des risques algorithmiques dans les administrations publiques.

Le premier chapitre a posé le contexte général en identifiant les risques systémiques liés à l'utilisation croissante des systèmes algorithmiques dans le secteur public, et a défini la méthodologie itérative adoptée sur 12 semaines. Le deuxième chapitre a formalisé le cadre conceptuel et méthodologique en établissant les formules mathématiques des trois modules : normalisation et scoring pour O-PAGe, agrégation multi-niveaux et indices de conformité pour M-PAGe, matrice d'impact et calcul d'efficacité pour I-PAGe.

Le troisième chapitre a détaillé l'analyse et la conception à travers l'identification de sept acteurs, la spécification des cas d'utilisation, les diagrammes de classes et la définition d'une architecture technique distribuée (React, Django REST Framework, PostgreSQL, Docker Compose). Le quatrième chapitre a présenté la réalisation concrète de l'application avec des interfaces différenciées par rôle, des tableaux de bord analytiques, des campagnes d'évaluation structurées, un module de simulation avec export, et des analyses comparatives multi-projets.

La plateforme livrée démontre la faisabilité d'une approche unifiée de la gouvernance algorithmique, intégrant observation, maturité et simulation d'impact au sein d'un même écosystème technique.

\vspace{0.3cm}
\noindent\textbf{Perspectives:} En vue d'un déploiement réel en environnement de production, la plateforme pourrait bénéficier de l'adoption de pratiques DevOps permettant d'assurer la fiabilité, la traçabilité et l'évolutivité continue de la solution.
